% ======================================================================
% Ontology of Continua -- Core 1.3
% Cross-K module: crossk_k8_k9.tex
% Cross-level relations between K8 and K9
% ======================================================================

\subsubsection{Межуровневая структура K8–K9}
\label{sec:crossk-k8-k9}

Переход $K_8 \rightarrow K_9$ обозначает возникновение
\textbf{теоретического / научного / метапарадигмального континуума} из
цивилизационно-технологического континуума. Если $K_8$ поддерживает
системы популяционного уровня, инфраструктуры, символическую кодировку и
технологические циклы, то $K_9$ вводит:
\begin{itemize}
  \item формальные теоретические рамки,
  \item логические и методологические оси,
  \item пороги согласованности, доказательности и объяснительной силы,
  \item потоки аргументации, доказательства и критики,
  \item парадигмальные циклы и научная эволюция,
  \item рекурсивные структуры, способные описывать нижние уровни и сами себя.
\end{itemize}

$K_9$ — первый уровень, где знание становится явной структурной
сущностью со своими условиями непрерывности.

% ----------------------------------------------------------------------
\subsubsection{1. Задействованные уровни и их роли}

\paragraph{K8 (цивилизационно-технологический континуум).}
$K_8$ предоставляет:
\begin{itemize}
  \item символические системы, письмо, математику, формальные языки,
  \item институты передачи и накопления знаний,
  \item информационные сети и кодированные репозитории,
  \item вычислительные и технологические возможности,
  \item пороги $\Theta_8$ (энергия, логистика, информационная когерентность).
\end{itemize}

Однако $K_8$ не включает:
\begin{itemize}
  \item формальные мета-структуры сравнения теорий,
  \item логические пороги согласованности или противоречивости,
  \item явные показатели объяснительной или прогностической силы,
  \item структурированные циклы научной эволюции.
\end{itemize}

\paragraph{K9 (теоретический континуум).}
$K_9$ вводит:
\begin{itemize}
  \item оси $A^9$ (логика, онтология, методология, интерпретация),
  \item потенциалы $P^9$ (объяснительная сила, прогностическая точность,
        когерентность, парсимоничность),
  \item пороги $\Theta^9$ (логическая согласованность, эмпирическая адекватность,
        глубина эвристики, гёделевские ограничения),
  \item потоки $J^9$ (аргументы, доказательства, критика, сдвиги парадигм),
  \item циклы $C^9$ (нормальная наука, накопление аномалий, революция),
  \item меру $k(K_9)$, отражающую глобальную теоретическую согласованность.
\end{itemize}

Таким образом, $K_9$ — это метапредставительный слой над цивилизационной осью.

% ----------------------------------------------------------------------
\subsubsection{2. Совместные и наследуемые оси и пороги}

\paragraph{Наследование символики.}
Символические инфраструктуры $K_8$ становятся подложкой для теоретических
различий:
\[
A_{\mathrm{symbol}}(K_8) \rightarrow A^9.
\]
Письмо → математика → логика → формальные модели.

\paragraph{Новые оси.}
$K_9$ вводит:
\begin{itemize}
  \item логические оси (валидность, правила вывода),
  \item онтологические оси (типы сущностей, структурные предпосылки),
  \item методологические оси (модели, эвристики, парадигмы),
  \item интерпретационные оси (конкурирующие прочтения одной теории).
\end{itemize}

\paragraph{Новые пороги.}
$\Theta^9$ включает:
\begin{itemize}
  \item \textbf{порог логической согласованности} — его нарушение уничтожает
        теорию,
  \item \textbf{эмпирический порог} — недостаточная сопоставимость с данными,
  \item \textbf{эвристический порог} — недостаточная объяснительная глубина,
  \item \textbf{парадигменный порог} — несопоставимость с существующими рамками,
  \item \textbf{гёделевский порог} — ограничения неполноты.
\end{itemize}

Если в $K_8$ нарушается информационная когерентность $\Theta_I$:
\[
\Theta_I^{(K_8)\mathrm{death}} \Rightarrow \Omega(K_9)=\varnothing.
\]

% ----------------------------------------------------------------------
\subsubsection{3. Межуровневые потоки, циклы и напряжения}

\paragraph{Потоки.}
Теоретические потоки $J^9$ расширяют символические потоки $K_8$:
\[
J^9 = (J_{\mathrm{argument}},\ J_{\mathrm{proof}},\
       J_{\mathrm{critique}},\ J_{\mathrm{reinterpret}}).
\]

Эти потоки изменяют $A^9$, $P^9$ и границы $\partial\Omega(K_9)$.

\paragraph{Циклы.}
Научные циклы $C^9$ включают:
\[
C_{\mathrm{science}}:\
\text{theory} \rightarrow \text{prediction} \rightarrow \text{experiment}
\rightarrow \text{anomaly} \rightarrow \text{revision} \rightarrow \text{theory}.
\]

Теоретический континуум существует только тогда, когда хотя бы один такой
стабильный цикл сохраняется.

\paragraph{Напряжение.}
Межуровневое напряжение:
\[
T_{8,9} =
f(\Theta^9 - P^9,\ A^9 - A_{\mathrm{symbol}},\
  J^9 - J_{\mathrm{info}},\ \Theta_8 - \Theta^9).
\]
Если $T_{8,9}$ превышает размерностный порог, парадигмы коллапсируют в
нелогические символические структуры.

% ----------------------------------------------------------------------
\subsubsection{4. Условия рождения и исчезновения между уровнями}

\paragraph{Рождение \texorpdfstring{$K_9$}{K_9}.}
Теоретический континуум возникает, когда:
\[
T_8 > \Theta_{8,\mathrm{dim}}
\quad\text{и}\quad
A^9 \in M_9 \setminus A(K_8),
\]
т.е. появляются новые логические/онтологические оси, которые нельзя свести к
технологии или цивилизационной организации.

Расширение:
\[
\Omega(K_9) = \Omega(K_8) \cup \Delta\Omega_{\mathrm{theory}}.
\]

\paragraph{Распространение гибели.}
Если $\Theta_8$ нарушена (цивилизационный коллапс), $K_9$ теряет субстрат.

Если $\Theta^9$ нарушена (логическая несогласованность, эмпирическое
опровержение, коллапс парадигмы), теоретический континуум разрушается, но
$K_8$ остаётся:
\[
\Omega(K_9) \rightarrow \varnothing.
\]

% ----------------------------------------------------------------------
\subsubsection{5. Концептуальные примеры}

\paragraph{Формализация знаний.}
Символические технические системы $K_8$ позволяют математические и логические
структуры, формирующие оси $K_9$.

\paragraph{Цикл научной революции.}
По мере накопления аномалий сверх $\Theta^9_{\mathrm{heuristic}}$:
\[
T_{8,9} > \Theta^9_{\mathrm{paradigm}},
\]
происходит смена парадигмы, создающая новую область $\Omega(K_9)$.

\paragraph{Случай коллапса.}
Если информационная когерентность в $K_8$ падает ниже $\Theta_I$, теоретическое
производство становится невозможным — признак цивилизационного коллапса.

Следовательно, переход K8→K9 можно сжато охарактеризовать как возникновение
формальных теорий, логики и научной эволюции из технологических,
символических и институциональных инфраструктур.
